\documentclass[a4paper, 11pt]{article}
\usepackage{comment}
\usepackage{lipsum} 
\usepackage{fullpage} %cambiar margen
\usepackage[a4paper, total={7in, 10in}]{geometry}

\usepackage{amssymb,amsthm} 
\usepackage{amsmath}
\newtheorem{theorem}{Theorem}
\newtheorem{corollary}{Corollary}
\usepackage{graphicx}
\usepackage{tikz}
\usetikzlibrary{arrows}
\usepackage{verbatim}
%\usepackage[numbered]{mcode}
\usepackage{float}
\usepackage{tikz}
\usetikzlibrary{shapes,arrows}
\usetikzlibrary{arrows,calc,positioning}
\usepackage{mathpazo} %tipo de letra 
\usepackage[utf8]{inputenc} %codificación
\usepackage[T1]{fontenc} %digitación de tildes y ñ
\usepackage[spanish]{babel} %paquete de soporte español

\tikzset{
	block/.style = {draw, rectangle,
		minimum height=1cm,
		minimum width=1.5cm},
	input/.style = {coordinate,node distance=1cm},
	output/.style = {coordinate,node distance=4cm},
	arrow/.style={draw, -latex,node distance=2cm},
	pinstyle/.style = {pin edge={latex-, black,node distance=2cm}},
	sum/.style = {draw, circle, node distance=1cm},
}
\usepackage{xcolor}
\usepackage{mdframed}
\usepackage[shortlabels]{enumitem}
\usepackage{indentfirst}
\usepackage{hyperref}

\usepackage{listings}
\lstset{literate=
  {á}{{\'a}}1
  {é}{{\'e}}1
  {í}{{\'i}}1
  {ó}{{\'o}}1
  {ú}{{\'u}}1
  {Á}{{\'A}}1
  {É}{{\'E}}1
  {Í}{{\'I}}1
  {Ó}{{\'O}}1
  {Ú}{{\'U}}1
  {ñ}{{\~n}}1
  {ü}{{\"u}}1
  {Ü}{{\"U}}1
}

\lstdefinestyle{customc}{
  belowcaptionskip=1\baselineskip,
  breaklines=true,
  frame=L,
  xleftmargin=\parindent,
  language=Python,
  showstringspaces=false,
  basicstyle=\footnotesize\ttfamily,
  keywordstyle=\bfseries\color{green!40!black},
  commentstyle=\itshape\color{purple!40!black},
  identifierstyle=\color{blue},
  stringstyle=\color{orange},
}

\lstdefinestyle{customasm}{
  belowcaptionskip=1\baselineskip,
  frame=L,
  xleftmargin=\parindent,
  language=[x86masm]Assembler,
  basicstyle=\footnotesize\ttfamily,
  commentstyle=\itshape\color{purple!40!black},
}

\lstset{escapechar=@,style=customc}



\renewcommand{\thesubsection}{\thesection.\alph{subsection}}

\newenvironment{problem}[2][Ejercicio]
{ \begin{mdframed}[backgroundcolor= red!50] \textbf{#1 #2} \\}
	{  \end{mdframed}}

% Define solution environment
\newenvironment{solution}
{\textcolor{blue}{\textbf{\textit{Solución:\\\noindent}}}}


\renewcommand{\qed}{\quad\qedsymbol}

% \\	
\begin{document}
	\noindent
	%%%%%%%%%%%%%%%%%%%%%%%%%%%%%%%%%%%%
	
	\begin{minipage}[b][1.2cm][t]{0.8\textwidth}
		\large\textbf{César Isaí García Cornejo} \hfill \textbf{Examen Parcial}  \\
		cesar.cornejo@cimat.mx \hfill \\
		\normalsize Series de Tiempo \hfill Semestre 3\\
	\end{minipage}
	
	\hspace{14.4cm}
	\begin{minipage}[b][0.03cm][t]{0.12\linewidth}
		
		\vspace{-2.2cm}
		%%%La Ruta dependerá de donde este alojado el main y la imagen
		\includegraphics[scale=0.3]{Figures/EscudoCimat.png}
	\end{minipage}
	
	\noindent\rule{7in}{2.8pt}
	
	%%%%%%%%%%%%%%%%%%%%%
	%%%%%%%%%%%%%%%%%%%%%%%%%%%%%%%%%%%%%%%%%%%%%%%%%%%%%%%%%%%%%%%%%%%%%%%%%%%%%%%%%%%%%%%%%%%%%%%%%%%%%%%%%%%%%%%%%%%
	% Problem 1
	%%%%%%%%%%%%%%%%%%%%%%%%%%%%%%%%%%%%%%%%%%%%%%%%%%%%%%%%%%%%%%%%%%%%%%%%%%%%%%%%%%%%%%%%%%%%%%%%%%%%%%%%%%%%%%%%%%%%%%%%%%%%%%%%%%%%%%%%
	\setlength{\parskip}{\medskipamount}
	\setlength{\parindent}{0pt}
 %///////////////////////////////////////////////////



\begin{problem}{1} 
    Considera la ST
    \begin{align*}
        Y_t = \text{0.4} Y_{t-1} + \text{0.45}Y_{t-2} + Z_t + Z_{t-1} + \text{0.25}Z_{t-2}
    \end{align*}
    donde $Z_t \sim WN(0,\sigma^2)$.
    \begin{enumerate}
        \item Expresa esta ecuación en términos del operador ``backshift'' $B$; es decir, escríbela como una ecuación en $B$, y determina el orden $(p,d,q)$ de este modelo.
        \item ?` Puedes simplificar esta ecuación? ?` Cuál es el orden después de la simplificación?
        \item determina si este modelo es causal y/o invertible.
        \item Si el modelo es causal, encuentra la forma general de los coeficientes $\psi_j$'s tal que $$Y_t = \sum_{j = 0}^{\infty} \psi_j Z_{t-j}$$.
        \item Si el modelo es invertible, encuentra la forma general de los coeficientes $\pi_j $ tal que $$Z_t = \sum_{j=0}^{\infty} \pi_j Y_{t-j}$$. 
    \end{enumerate}
\end{problem}

\begin{solution} 
  \begin{enumerate}
    \item Expresamos la serie de tiempo en la forma de un ARIMA(p,d,q)
    \begin{align*}
      Y_t -\text{0.4}Y_{t-1} - \text{0.45} Y_{t-2} &= Z_t + Z_{t-1} + \text{0.25} Z_{t-2},\\
      \left(1-\text{0.4}  B -\text{0.45} B^2  \right) Y_t &= \left(1+B+\text{0.25}B^2 \right) Z_t ,\\
      \left(1 + (\text{0.9} - \text{0.5} )B + (\text{0.5}  ) (-\text{0.9} )B^2\right)  Y_t &= \left(1 + 2(\text{0.5} ) B + (\text{0.5})^2 B^2\right) Z_t ,\\
      (1-\text{0.9} B )(1+\text{0.5} B) Y_t &= (1+\text{0.5} B)^2 Z_t,\\
      \phi'(B) Y_t &= \theta'(B) Z_t. 
    \end{align*}
    donde $\phi'(B) = (1-\text{0.9} B )(1+\text{0.5} B)$ y $\theta' (B) = (1+\text{0.5} B)^2$ como no hay raíces unitarias, el valor de $d = 0$. Entonces, esta serie de tiempo es un ARIMA(2,0,2).
    
    \item Recordemos que para evitar tener problemas de identificabilidad es necesario simplificar los factores, de ser posible. En este caso tenemos que la simplificación es
    \begin{align*}
      (1-\text{0.9} B )Y_t &= (1+\text{0.5} B) Z_t,\\
      \phi(B) Y_t &= \theta(B) Z_t. 
    \end{align*}
    donde $\phi(B) = 1-\text{0.9} B $ y $\theta(B) = 1+\text{0.5} B$ y el orden de esta serie de tiempo es (1,0,1).
    \item Recordemos que un modelo es causal si y solo si todas las raíces de $\phi (z)$ son mayores que uno en norma. Un modelo es invertible si y solo si las raíces de $\theta(z) $ son mayores que uno en norma. Así
    para $\phi(z)$
    \begin{align*}
      0 = \phi(z) = 1-\text{0.9} z    \:\:\: \Rightarrow \:\:\: z_1 = 10/9
    \end{align*}
    como $|z_1| > 1$ entonces la ST es causal. Luego,
    \begin{align*}
      0 = \theta(z) = 1 + \text{0.5} z  \:\:\: \Rightarrow \:\:\: z_2 = -2
    \end{align*}
    como $|z_2| > 1 $ entonces la ST es invertible.
    \item Ahora que ya probamos que la ST es causal e invertible, para obtener la expresión solicitada hacemos la siguiente manipulación
    \begin{align*}
      Y_t &= \frac{\theta(B)}{\phi(B)} Z_t ,\\  
      Y_t &= \left(1+\text{0.5}B \right) \frac{1}{1-\text{0.9} B} Z_t
    \end{align*}
    por expansión en serie geométrica
    \begin{align*}
      Y_t &= (1+ \text{0.5} B ) \left(1 + \text{0.9}  B + \text{0.9}^2B^2 + \text{0.9}^3 B^3 + \cdots \right) Z_t,\\
      &= \left(1 + (\text{0.9})^0(\text{0.9} + \text{0.5} )B + (\text{0.9} ^2 + \text{0.9} \cdot \text{0.5} ) B^2 + (\text{0.9} ^3 + \text{0.9} ^2 \cdot \text{0.5} )B^3 + \cdots\right) Z_t,\\
      &= \left(1 + \sum_{j=1}^{\infty}\text{1.4} \cdot (\text{0.9} )^{j-1}B_j   \right) Z_t ,\\
      &= Z_t + \sum_{j=1 }^{\infty} \text{1.4} (\text{0.9} )^{j-1} Z_{t-j} 
    \end{align*}
    Comparando con la expansión causal del enunciado, $\psi_0 = 1$ y $\psi_j = \text{1.4} (\text{0.9} )^{j-1}$ para $j = 1,2,\cdots$.
    \item Semejante al anterior tomamos la representación
    \begin{align*}
      Z_t &= \frac{\phi(B)}{\theta(B)}Y_t 
    \end{align*}
    con la expansión de la serie geométrica y simplificando
    \begin{align*}
      Z_t &= \frac{1 - \text{0.9} B }{1+ \text{0.5} B} Y_t ,\\
      &= (1-\text{0.9} B)(1-\text{0.5} B + \text{0.5} ^2 B^2  - \text{0.5} ^3 B^3 + \cdots)Y_t ,\\
      &= \left(1 + (-\text{0.5} - \text{0.9} )B + (\text{0.5}^2 + (\text{-0.5})(-\text{0.9} ) )B^2 +  (-\text{0.5} ^3 -(\text{0.9} )(\text{-0.5} ^2)) B^3 + \cdots \right) Y_t,\\
      &= \left(1 + \sum_{j=1}^{\infty} \left((-\text{0.5} )^{i} - \text{0.9} (-\text{0.5} )^{i-1}\right) \right) Y_t,\\
      &= \left(1 + \sum_{j=1}^{\infty}(-\text{0.5} )^{i-1}(-\text{0.5} - \text{0.9} )\right) Y_t,\\
      &= \left(1 + \sum_{j=1}^{\infty}(-\text{0.5})^{i-1}(-\text{1.4})\right) Y_t,
    \end{align*}
    Comparando con la expresión invertible dada por el enunciado. Tenemos, $\pi_0 = 1$ y $\pi_j = -\text{1.4} (-\text{0.5} )^{j-1}$ para $j = 1,2, \cdots$
  \end{enumerate}
    
\end{solution}

\begin{problem}{2} 
    Asume que una ST satisface la ecuación en diferencias
    \begin{align*}
        Y_t = \left\{\begin{matrix}
            1 & t = 0,1 \\ 
            \text{0.4} Y_{t-1}  +\text{0.77} Y_{t-2} + e_t &\:\:\:\:\:\:\: t=2,3, \cdots 
            \end{matrix}\right.
    \end{align*}
    donde las $e_t$'s son v.a. no-correlacionadas(0,1). Sea $(Y_2,Y_4,Y_5) =$ (2,3,2.9).
    \begin{enumerate}
        \item Predice $Y_6$ y $Y_7$. Calcula la matriz de covarianzas de tu predicción. \textit{Hint: Usa el mejor predictor general} 
        \item ?`Cuál es el límite $\lim_{s\rightarrow\infty}\mathbb{E}\left [\left(Y_{5+s}- \hat{Y_{5+s}} \right)^{2}\right ]$ ?
    \end{enumerate}
\end{problem}

\begin{solution} 
  

\end{solution}

\begin{problem}{3} 
    Sea $Y_t = Z_t - \theta Z_{t-1}$, $Z_t \sim WN(0,\sigma^2)$.
    \begin{enumerate}
        \item Calcula la función de correlación $\rho(k)$ de $Y_t$.
        \item Supón que $\rho(1) = $ 0.4. ?` Cuál preferirías? Proporciona una explicación corta. 
        \item En lugar de un modelo MA(1), asume que $Y_t$ es un modelo MA($\infty$) dado por 
        \begin{align*}
            Y_t = Z_t + C (Z_{t-1} + Z_{t-2 } + \cdots),
        \end{align*}
        \item Si $\{Y_t \}$ en la ecuación anterior se diferencía (i.e.,$ X_t = Y_t - Y_{t-1}$), muestra que $X_t$ es un modelo estacionario MA(1).
        \item Encuentra la función de autocorrelación de $\{X_t\}$.
    \end{enumerate}
\end{problem}

\begin{solution} 
  \begin{enumerate}
    \item Calculemos primeramente la función de autocovarianza
    \begin{align*}
      \gamma(k) = Cov(Y_t,Y_{t+k}).
    \end{align*}
    Para $k = 0$ 
    \begin{align*}
      \gamma(0) &= \mathbb{V}ar[Y_t ],\\
      &= \mathbb{V}ar[Z_t - \theta Z_{t-1}] ,\\
      &= \mathbb{V}ar[Z_t] + \theta^2 \mathbb{V}ar[Z_{t-1}] 
    \end{align*}
    donde se usó la no-correlación del ruido blanco $Z_t$. Luego,
    \begin{align*}
      \gamma(0) = (1+ \theta^2)\sigma^2
    \end{align*}
    Para $k = 1$
    \begin{align*}
      \gamma(1) &= Cov(Y_t, Y_{t+1}) ,\\
      &= Cov(Z_t - \theta Z_{t-1}, Z_{t+1} - \theta Z_{t} ),\\
      &= -\theta \sigma^2
    \end{align*}
    Lo mismo es para $k= -1$. Para $|k| \geq 2$ no hay términos con covarianza no nula, es decir su covarianza es nula. Resumiendo
    \begin{align*}
      \gamma(k) = \left\{\begin{matrix}
        (1+ \theta^2) \sigma^2&, k = 0 \\ 
        -\theta \sigma^2 &, |k| = 1\\
        0 &, |k| \geq 2 
      \end{matrix}\right.
    \end{align*}
    Se sigue que
    \begin{align}
      \rho(k) = \frac{\gamma(k)}{\gamma(0)} = \left\{\begin{matrix}
        1 & ,k = 0 \\ 
        \frac{-\theta}{1 + \theta^2} &, |k| = 1\\
        0 &, |k| \geq 2 
      \end{matrix}\right.
      \label{3.01}
    \end{align}
    es la función de autocorrelación.
    
    \item Supongamos que $\rho_1 = \rho(1) =$ 0.4. De (\ref{3.01}) 
    \begin{align*}
      \rho_1 = -\frac{\theta}{1+ \theta^2},\\
      \rho_1 + \rho_1 \theta^2 + \theta = -\theta,\\
      \rho_1 \theta^2 +\theta + \rho_1 = 0.
    \end{align*}
    Resolviendo la ecuación cuadrática
    \begin{align*}
      \theta = \frac{-1\pm \sqrt{1-4\rho_1^2}}{2\rho_1} = \frac{-1 \pm \text{0.6} }{\text{0.8} }
    \end{align*}
    obteniendo las raíces
    \begin{align*}
      \theta_+ &= -0.5,\\
      \theta_- &= -2.
    \end{align*}
    Nos quedamos con $\theta = $-0.5 porque con está tenemos las raíces del polinomio característico mayores que uno en norma. Lo que implica que la ST es invertible.
    
    \item Sabemos que una serie de tiempo es estacionario si tanto su media como su función de covarianza no dependen del tiempo $t$ (si la función de covarianza es finita). Analicemos dichas características. Para la esperanza 
    \begin{align*}
      \mathbb{E}\left [Y_t \right ] &= \mathbb{E}\left [ Z_t + C(Z_{t-1} + Z_{t-2} + \cdots )\right ],\\
      &= 0
    \end{align*}
    por linealidad de la esperanza. 

    Consideremos la varianza de la ST
    \begin{align*}
      \mathbb{V}ar[Y_t  ] &= \mathbb{V}ar[Z_t + C(Z_{t-1} + Z_{t-2} + \cdots )],\\
      &= \sigma^2 + C^2 (\sigma^2 + \sigma^2 + \cdots)
    \end{align*}
    donde se usó la no-correlación del ruido blanco. Si $C \neq 0$ entonces la varianza es infinita lo que implica que no puede ser estacionaria ya que la definición tiene como hipótesis varianza finita.

    \item Haciendo el calculo explícito
    \begin{align*}
      X_t &= Y_t - Y_{t-1},\\
      = &Z_t + C(Z_{t-1} + Z_{t-2} + Z_{t-3} + \cdots) +\\
      & -Z_{t-1} - C(Z_{t-2} + Z_{t-3} + Z_{t-3} + \cdots )\\
      &= Z_t - Z_{t-1} + C(Z_{t-1})
    \end{align*}
    lo que se simplifica al MA(1) 
    \begin{align*}
      X_t = Z_t + (C-1)Z_{t-1}
    \end{align*}
    que sabemos que es estacionario.

    \item Como es un MA(1) y ya calculamos la correlación en (\ref{3.01}) con $\theta = C-1$. Tenemos que la función de autocorrelación es
    \begin{align*}
      \rho(k) = \left\{\begin{matrix}
        1 & ,k = 0 \\ 
        -\frac{(C-1)}{1 + (C-1)^2} &, |k| = 1\\
        0 &, |k| \geq 2 
        \end{matrix}\right.
    \end{align*}
    lo que concluye el ejercicio. 
    
    
    
    
    
    
    
    
    
    
    
    
    
    
    
    
  \end{enumerate}

\end{solution}





\end{document}